\chapter{Host Organization and Project Context}
\label{chap:host}

\section{Introduction}

This chapter introduces the organization in which the project was carried out, describes the context that led to the project, states the problem addressed, and lists the objectives that guided the rest of the work. It closes with a short summary that leads into Chapter~\ref{chap:sota}.

\section{Host organization}

\todofill{Company name}, described here, hosted this final year project.

\subsection{Presentation}

\todofill{Add a short (half-page) presentation of the host company: sector of activity, size, main products or services, location. If this was an independent / self-directed project rather than a company internship, replace this subsection with a short paragraph stating that explicitly, and describe instead the academic framework under which the project was carried out.}

\subsection{Department}

\todofill{Name the department, service, or team in which the project took place, and its role inside the organization.}

\section{Project context}

The project was carried out in the broader context of automotive cybersecurity. Modern vehicles, and electric vehicles in particular, are built around internal networks of \gls{ecu}s that communicate over a \gls{can} bus. This bus was designed for reliability, not for security, and it has been repeatedly shown in the literature to be vulnerable to message injection, flooding, and replay attacks (see Chapter~\ref{chap:sota}).

The project, \textbf{CANvisionNative}, is a hybrid-architecture \gls{ids} for vehicle \gls{can} bus networks. It combines:

\begin{itemize}
    \item a Windows desktop client (\gls{wpf}, \gls{mvvm}) for visualization, replay control, and analyst interaction;
    \item a Python backend (FastAPI) that decodes \gls{can} frames, extracts features, scores anomalies, and generates natural-language explanations;
    \item a machine-learning layer trained on more than seven million \gls{can} frames across eighteen vehicle profiles;
    \item support for three operating modes: replay of recorded logs, live simulation, and live hardware capture.
\end{itemize}

The system was designed to be usable both as a forensic tool — importing a recorded log and reviewing detected anomalies after the fact — and, once hardware is connected, as a live monitoring tool.

\section{Problem statement}

Three problems motivated this project.

\textbf{First}, the \gls{can} bus has no built-in security. Any device connected to the bus can send frames, and no frame carries a sender identity or a signature. An intrusion detection system is therefore one of the only practical ways to add a security layer to an existing vehicle without redesigning its electrical architecture.

\textbf{Second}, a single detection strategy is not enough. Timing-based detection catches flooding attacks but misses slow, low-rate attacks. Payload-based detection catches out-of-range values but misses attacks that stay within plausible ranges. Machine-learning models trained on one vehicle rarely transfer cleanly to another vehicle, because message layout, frequency, and normal-value ranges are vehicle-specific. This suggested a design combining several independent detection layers, fused into a single score, rather than relying on one algorithm.

\textbf{Third}, a raw anomaly score is not enough for a human analyst to act on. A list of alerts with numeric scores does not explain what happened, why the system flagged it, or what to do next. This motivated an explanation layer that turns a detected anomaly into a short, readable description of the likely attack pattern and a recommended action.

Beyond building the system, this project also required a rigorous validation phase, since a detection system that produces incorrect or misleading alerts is not trustworthy. This is documented in full in Chapter~\ref{chap:validation}, which describes four confirmed software defects that were found through evidence-based testing on a real one-hour vehicle recording, and their correction.

\section{Objectives}

The objectives of the project were:

\begin{enumerate}
    \item Design and implement a multi-layer \gls{can} bus intrusion detection pipeline combining timing, payload, machine-learning, and temporal-persistence signals into a single fusion score.
    \item Support multiple input sources: recorded logs (\gls{csv}, \gls{mf4}, candump, Vector \gls{can}), a live attack simulator, and live hardware capture through an ELM327 or SocketCAN adapter.
    \item Train and integrate vehicle-specific machine-learning models so that detection adapts to the vehicle being analysed instead of relying on one generic model.
    \item Build a desktop client that lets an analyst import a log, replay it, review detected anomalies, and export a report.
    \item Add a natural-language explanation layer that describes each anomaly and recommends an action, using a combination of rule-based reasoning and a Large Language Model.
    \item Validate the complete pipeline against a real vehicle recording, using measurable, reproducible evidence, and correct any defect found before considering the system ready.
\end{enumerate}

\section{Conclusion}

This chapter presented the organization and context in which CANvisionNative was developed, the problems that motivated it, and the objectives set at the start of the project. The next chapter reviews the state of the art on \gls{can} bus security, existing intrusion detection approaches, and the use of machine learning for this problem, in order to position the choices made in this project.
